%% kapa: exam=1 round=23 year=2012 period=2 form=A subject=accounting subject_ko=회계학 seq=20 number=20
%% source: Q-Net 감정평가사 기출문제 (cst00309) artlSeq=5052580 file=raw/1008/5052580/2012년 제23회 감정평가사 제1차시험 기출문제(2교시 A형).pdf
%% conversion: kordoc 4.24.51 -> tools/gen1.py (auto); review=proofread-sonnet; flags=none
\begin{question}{20}
무형자산에 관한 설명으로 옳지 않은 것은?
\begin{choices}
  \item 무형자산은 물리적 실체는 없지만 식별가능한 비화폐성자산을 말하며, 식별가능성은 자산의 분리가능성 여부와 계약상 또는 법적 권리여부에 의해 판단할 수 있다.
  \item 무형자산으로 분류되기 위해서는, 식별가능해야 하고, 통제하고 있어야 하며, 그로부터 미래경제적 효익이 창출되어야 한다.
  \item 광물자원과 관련된 지출은 해당자산이 기술적 실현가능성과 상업화가능성이 있다고 판단되어 개발 활동을 시작하면 이 단계부터의 지출을 탐사평가 자산이 아닌 개발비로 인식한다.
  \item 무형자산이 한정할 수 없는 기간동안 순현금유입을 창출할 수 있을 것으로 기대된다면 그 무형자산은 비한정내용연수를 가지며, 상각하지 아니한다.
  \item 무형자산의 공정가치는 합리적인 공정가치측정치의 범위의 편차가 자산가치에 비하여 유의적일 때 신뢰성 있게 측정할 수 있다.
\end{choices}
\end{question}
\answer{5}
